\documentclass{article}
\usepackage{amsmath,amssymb,amsthm}
\usepackage{hyperref}
\usepackage{tcolorbox}
\tcbuselibrary{breakable}

\newtcolorbox{accept}[1][]{colback=green!5,colframe=green!60!black,
  fonttitle=\bfseries,title={Accepted: #1},breakable}
\newtcolorbox{partial}[1][]{colback=yellow!5,colframe=orange!60!black,
  fonttitle=\bfseries,title={Partially Accepted: #1},breakable}
\newtcolorbox{reject}[1][]{colback=red!5,colframe=red!60!black,
  fonttitle=\bfseries,title={Not Accepted (with reason): #1},breakable}

\title{Response to Reviewer Gem31}
\date{April 2026}
\begin{document}
\maketitle

\section{Overall Assessment of Review}

Gem31 provides a positive, constructive review that correctly identifies two
genuine gaps and proposes mathematically sound fixes. The review is shorter
than the other three but focuses on the two most practically impactful issues.

\section{Gap 1: Metric Approximation Assumption A4'}

\begin{accept}[Sobolev interpolation path -- accepted in principle]
Gem31 correctly identifies that the finite-difference argument gives
$\|G^* - g\|_{\mathrm{op}} = O(\delta_{\mathrm{rec}}/r_n + r_n)$ which blows
up as $n \to \infty$ for fixed $\delta_{\mathrm{rec}}$.

The proposed Gagliardo-Nirenberg / Sobolev interpolation approach is sound in
spirit. However, the correct and simpler resolution (also provided by reviewers
Op46MX and Snn46) is:

\textbf{The capacity condition $\delta_{\mathrm{rec}} = O(r_n^2)$ suffices.}
Under this condition, $\delta_{\mathrm{rec}}/r_n = O(r_n) \to 0$, and the
metric bound becomes $\|G^* - g\|_{\mathrm{op}} = O(r_n)$, which is exactly
the form needed for Theorem 2 (absorbed into the $C_1 r_n$ term).

\textbf{Mathematical derivation (accepted from Snn46/Op46MX):}
By the finite-difference formula, $\|J_{f^*}(z_i) - J_{f_{\mathrm{true}}}(z_i)\|_{\mathrm{op}}
\leq 2\delta_{\mathrm{rec}}/r_n + C\kappa r_n$. Under $\delta_{\mathrm{rec}}
= O(r_n^2)$, this equals $O(r_n)$. Then $\|G^* - g\|_{\mathrm{op}} =
O(\|J^* - J_{\mathrm{true}}\|_{\mathrm{op}}) = O(r_n)$.

\textbf{Action}: Add a Proposition/Remark after Assumption A4' in the paper
with this derivation. State explicitly that A4' holds whenever
$\delta_{\mathrm{rec}} = O(r_n^2)$ (which holds for decoders with sufficient
capacity).
\end{accept}

\begin{partial}[Sobolev norm regularization -- partially accepted]
Gem31 proposes adding explicit Sobolev norm regularization to the loss to control
the Jacobian error directly. This is a valid architectural improvement but is
not needed for the theoretical result: the capacity condition suffices.

\textbf{Decision}: We note this as a potential extension (the Riemannian loss
$\mathcal{L}_{\mathrm{Riem}}$ already implicitly regularizes the Jacobian by
training the edge decoder to predict $J_f(z_i)\Delta z_{ij}$). We will not
add explicit Jacobian regularization to the current paper but will mention it
as a future direction.
\end{partial}

\section{Gap 2: Topological Obstruction for the Torus}

\begin{accept}[Core diagnosis is correct -- accepted]
Gem31 correctly identifies that $\mathbb{R}^2$ is simply connected while
$T^2$ has fundamental group $\mathbb{Z} \times \mathbb{Z}$, so any continuous
map $\mathbb{R}^2 \to T^2$ must introduce singularities or folds. This is the
primary mathematical cause of the torus performance gap.

\textbf{Additional finding (from reviewers op46, Op46MX, Snn46):} The torus
experiment has a more severe problem -- the data is generated on the
\emph{standard embedded torus in $\mathbb{R}^3$} which has NONZERO Gaussian
curvature $K(\phi) = \cos\phi / [r(R+r\cos\phi)]$, while the evaluation uses
the \emph{flat torus geodesic formula}. This is a data-evaluation mismatch
that partially invalidates the torus results independently of the topology issue.
\end{accept}

\begin{accept}[Clifford torus fix -- accepted]
Gem31 proposes using the Clifford torus in $\mathbb{R}^4$:
$(R\cos\theta, R\sin\theta, r\cos\phi, r\sin\phi)$.

This is the correct fix:
\begin{itemize}
  \item Clifford torus is genuinely flat ($K = 0$ everywhere).
  \item The geodesic formula $d = \sqrt{(R\Delta\theta)^2 + (r\Delta\phi)^2}$
    is EXACT for the Clifford embedding.
  \item Embedding matrix becomes $A \in \mathbb{R}^{G \times 4}$.
\end{itemize}

\textbf{Action}: Implement \texttt{flat\_torus\_clifford()} in
\texttt{data/synthetic.py} and rerun all torus experiments.
\end{accept}

\begin{partial}[Periodic latent space -- partially accepted]
Gem31 proposes replacing $\mathbb{R}^2$ with a periodic latent space
$[0, 2\pi)^2$ with von Mises prior. This is architecturally correct and would
fully resolve the topological obstruction.

\textbf{Decision}: This is an important extension but changes the model
architecture significantly. We will implement the Clifford torus experiment
(simpler fix) and add the periodic latent space as a \textbf{future work}
direction in the paper. The topological obstruction will be formally stated
as a theorem (using the Brouwer/Borsuk-Ulam argument from Snn46).
\end{partial}

\section{Gem31 Follow-up: Validation of Our Decisions}

Gem31 reviewed our response and confirmed both resolutions:

\begin{enumerate}
  \item \textbf{Gap 1 (A4'):} Gem31 calls our capacity-condition fix ``brilliant''
    and ``much cleaner'' than the Sobolev regularization they originally proposed.
    The resolution is indeed simpler: $\delta_{\mathrm{rec}} = O(r_n^2)$ directly
    collapses the problematic $1/r_n$ factor without any architectural changes.
    No further action needed.

  \item \textbf{Gap 2 (Torus):} Gem31 confirms:
    (a) The Clifford torus fix is ``the perfect experimental fix'' for testing $K=0$.
    (b) Deferring the periodic latent space to future work while adding a formal
    topological obstruction theorem is ``very reasonable scoping for a NeurIPS
    paper -- explaining a fundamental limitation mathematically is often just as
    valuable as engineering around it.''
    No further action needed.
\end{enumerate}

\section{Action Items Accepted from Gem31}

\begin{enumerate}
  \item Add capacity-condition derivation of A4': $\delta_{\mathrm{rec}} = O(r_n^2)$
    implies A4' holds and Theorem 2 becomes unconditional (absorbed into $C_1 r_n$).
  \item Switch torus experiment to Clifford embedding in $\mathbb{R}^4$.
  \item Add topological obstruction theorem (formal, from Snn46's bridge proof)
    to explain why the torus isometry floor is a fundamental limitation, not
    an optimization failure.
  \item Update the torus curvature baseline: since the Clifford torus is
    genuinely flat, the theoretical prediction $\hat{\kappa}^* = 0$ will be
    valid, and the experiment will cleanly test the curvature detection claim.
\end{enumerate}

\end{document}
